\section*{Bab \ref{ch:g6:pure-substances-and-mixtures} --- Zat Murni dan Campuran}

\begin{solution}{exo:g6:pure-substances-and-mixtures:1}
\omterm{def:g6:pure-substances-and-mixtures:pure-substance}{Zat murni}: air suling, batang besi murni, kristal gula. \omterm{def:g6:pure-substances-and-mixtures:mixture}{Campuran}:
\omterm{def:g4:air-a-mixture-of-gases:air}{udara}, air laut, susu.
\end{solution}

\begin{solution}{exo:g6:pure-substances-and-mixtures:2}
Homogen: air asin, air keran. Heterogen: granit, minyak di atas air,
muesli.
\end{solution}

\begin{solution}{exo:g6:pure-substances-and-mixtures:3}
Satu jenis zat tertentu, dengan nama dan sifatnya sendiri: misalnya
air, garam, etanol (juga gula, oksigen, besi).
\end{solution}

\begin{solution}{exo:g6:pure-substances-and-mixtures:4}
Suhu didih yang tetap menunjukkan \omterm{def:g6:pure-substances-and-mixtures:pure-substance}{zat murni}; \qty{78}{\celsius} adalah
suhu didih etanol.
\end{solution}

\begin{solution}{exo:g6:pure-substances-and-mixtures:5}
Jus itu mengandung banyak \omterm{def:g6:pure-substances-and-mixtures:species}{spesies kimia} (air, gula-gula, asam, vitamin,
zat perisa): jus itu \omterm{def:g6:pure-substances-and-mixtures:mixture}{campuran}. Kata ``murni'' pada kotaknya hanya
berarti bahwa tidak ada yang ditambahkan.
\end{solution}

\begin{solution}{exo:g6:pure-substances-and-mixtures:6}
Bukan: \omterm{def:g6:pure-substances-and-mixtures:pure-substance}{zat murni} mendidih pada suhu yang tetap. Cairan itu \omterm{def:g6:pure-substances-and-mixtures:mixture}{campuran},
kemungkinan air dengan sesuatu yang \omterm{def:g2:mixing-and-dissolving:dissolve}{larut} di dalamnya, misalnya air
asin.
\end{solution}

\begin{solution}{exo:g6:pure-substances-and-mixtures:7}
Saringan dapat memisahkan air berlumpur (butirannya tertahan di
kertas) dan jus berbulir (bulirnya tertahan). Saringan tidak dapat
memisahkan air asin (garamnya \omterm{def:g2:mixing-and-dissolving:dissolve}{larut}) ataupun minyak dari air (kedua
cairan itu lolos; keduanya dipisahkan dengan menuang lapisan atasnya).
\end{solution}

\begin{solution}{exo:g6:pure-substances-and-mixtures:8}
$78.7 \div 10 = \qty{7.87}{g}$ untuk \qty{1}{cm^3}: besi.
\end{solution}

\begin{solution}{exo:g6:pure-substances-and-mixtures:9}
$\frac{357}{478} \approx 0.747$: sekitar \qty{75}{\percent}.
\end{solution}

\begin{solution}{exo:g6:pure-substances-and-mixtures:10}
$2 \times 35 = \qty{70}{g}$ garam-garam. Persentase: $\frac{35}{1000} =
\qty{3.5}{\percent}$.
\end{solution}

\begin{solution}{exo:g6:pure-substances-and-mixtures:11}
Ukur suhu didih keduanya (masing-masing harus tetap selama mendidih, dan
keduanya harus sama), lalu timbang volume yang sama dari masing-masing
(massanya sama). Dua tetapan yang sama, dan keduanya tetap, menunjukkan
\omterm{def:g6:pure-substances-and-mixtures:pure-substance}{zat murni} yang sama.
\end{solution}

\begin{solution}{exo:g6:pure-substances-and-mixtures:12}
$\frac{10}{90 + 10} = \qty{10}{\percent}$ etanol. Cairan itu \omterm{def:g6:pure-substances-and-mixtures:mixture}{campuran}:
tidak memiliki suhu didih yang tetap dan tidak tercantum dalam tabel,
yang hanya memuat \omterm{def:g6:pure-substances-and-mixtures:pure-substance}{zat murni}.
\end{solution}

\begin{solution}{pb:g6:pure-substances-and-mixtures:1}
\textbf{1.} Tidak. \omterm{def:g6:pure-substances-and-mixtures:pure-substance}{Zat murni} dan \omterm{def:g6:pure-substances-and-mixtures:homogeneous}{campuran homogen} dapat tampak sama
persis: harus ada yang diukur.

\textbf{2.} Homogen: cairan itu jernih dan bagian-bagiannya tidak
terlihat.

\textbf{3.} Sebuah tetapan: suhu didih yang tetap selama cairan
mendidih (atau massa suatu volume yang diketahui, dibandingkan dengan
tabel).

\textbf{4.} A dan B, yang suhu didihnya tetap, adalah \omterm{def:g6:pure-substances-and-mixtures:pure-substance}{zat murni}. C
adalah \omterm{def:g6:pure-substances-and-mixtures:mixture}{campuran}.

\textbf{5.} A mendidih pada \qty{100}{\celsius}: air. B mendidih pada
\qty{78}{\celsius}: etanol.

\textbf{6.} A: $49.8 \div 50.0 = \qty{0.996}{g}$. B: $39.5 \div 50.0 =
\qty{0.790}{g}$. C: $51.2 \div 50.0 = \qty{1.024}{g}$.

\textbf{7.} Ya: sekitar \qty{1.0}{g} untuk air, \qty{0.79}{g} untuk
etanol.

\textbf{8.} Ketika airnya mendidih dan pergi, garamnya tertinggal,
sehingga cairan yang tersisa mengandung makin banyak garam; cairan itu
lalu mendidih pada suhu yang lebih tinggi.

\textbf{9.} Cairan C mengandung padatan terlarut: cairan itu larutan,
sebuah \omterm{def:g6:pure-substances-and-mixtures:mixture}{campuran}, seperti yang sudah ditunjukkan oleh cara mendidihnya.

\textbf{10.} $100 \times 1.024 = \qty{102.4}{g}$; $\frac{3.5}{102.4}
\approx 0.034$: sekitar \qty{3.4}{\percent} massanya adalah garam.

\textbf{11.} $\frac{35}{1000} = \qty{3.5}{\percent}$: sangat dekat
dengan hasil soal 10.

\textbf{12.} \qty{3.5}{g} dalam \qty{100}{mL}, jadi $3.5 \times 10 =
\qty{35}{g}$ garam per liter. Dengan keasinannya dan cara mendidihnya,
C adalah air laut; A air suling dan B etanol.
\end{solution}
